\section{Frequency Relationship}
\label{Chap:Frequency Relationship}
\begin{frame}[t]{Chapter \autoref{Chap:Frequency Relationship}}
\tableofcontents[sections = \value{section}]
\end{frame}

\begin{frame}[c]
% c for center

\begin{center}
\Huge \textcolor{green}{Frequency Relationship}
\end{center}

\end{frame}

\subsection{Analog and Discrete Frequencies}
\begin{frame}[t,allowframebreaks]{Analog and Discrete Frequency Relationship}

\begin{itemize}

\item $ -\infty < \, \omega_c \left(\text{rad/sec}\right) \, < \, + \infty \, \leftrightarrow$ analog frequency

\item $ - \pi < \, \omega_d \left(\text{rad}\right) \, < \, + \pi \, \leftrightarrow$ discrete frequency

\end{itemize}

\begin{equation}
\begin{array}{rcl}
\omega_c \left(\text{rad/sec}\right) \, &=& \, \omega_d \, \times \, F_s \\ \\
&=& \, \displaystyle\frac{\omega_d}{T_s}  \\ \\

\omega_d \, &=& \, \omega_c \left(\text{rad/sec}\right) \, \times \, T_s \\ \\
&=& \, \displaystyle\frac{\omega_c}{F_s}
\end{array}
\label{Eq:Freq_relation:analog_discrete_freq}
\end{equation}



\end{frame}


\subsection{DTFT and Ideal Impulse CTFT Relation}
\begin{frame}[t]{DTFT and Ideal Impulse CTFT Relation}

\begin{itemize}

\item Small Recall:

	\begin{itemize}
	\item When we sample a signal in time domain, the original spectrum will get replicated: $X(f_c) \, \rightarrow \, \text{time domain sampling} \rightarrow X_{\delta}(f_c)$
	
		\begin{itemize}
		\item $X_{\delta}(f_c)$ is replicated at multiple of $F_s$ 
		\item Don't forget that we have a scaling amplitude factor $= \, \displaystyle\frac{\text{Amp}}{T_s}$
		\end{itemize}
	
	\item Now if we have a discrete signal $x[n]$, its DTFT $X(f_d)$ is replicated at multiple of $2 \pi$
	
		\begin{itemize}
		\item However, don't forget we have a relation between Sampled CTFT $X_{\delta}(f_c)$ and the DTFT $X(f_d)$ using equation (\autoref{Eq:Freq_relation:analog_discrete_freq})
			
	
	\end{itemize}

\end{itemize}

\end{itemize}

\end{frame}
